% TEMPLATE-MILC-v3.tex - plantilla maestra UNICA de las guias MILC v3.
%
% La usan las tres familias de guias, cada una con su driver:
%   · Grados 6-11  -> scripts/build-guias-g11.py
%   · Bebras       -> scripts/build-guias-bebras.py
%   · Semillero    -> scripts/build-guias-semillero.py
%
% Antes habia tres copias casi identicas (~400 lineas iguales, 6 distintas):
% cada mejora habia que aplicarla tres veces, y en la practica no se hacia
% (el motor de recursos vivia solo en dos de ellas). Lo que varia entre
% familias esta parametrizado en 6 placeholders que cada driver llena
% (se nombran aqui SIN su sintaxis real: el builder reemplaza por texto plano,
% tambien dentro de los comentarios, y un valor multilinea rompe el preambulo):
%
%   HEADER_IZQ        encabezado izquierdo de pagina
%   HEADER_DER        encabezado derecho de pagina
%   PORTADA_ETIQUETA  rotulo sobre el numero grande (GRADO/MOMENTO/MODULO)
%   PORTADA_SERIE     linea de serie de la portada
%   PORTADA_ID        identificador de la guia en la portada
%   PORTADA_CIRCULO   el circulito de grado (°), o vacio si no aplica
%
% El resto de claves las documenta content/guias/_SCHEMA.md. El script valida
% que no queden placeholders sin reemplazar antes de escribir el archivo final.

\documentclass[11pt,letterpaper]{article}
\usepackage[margin=1.75cm]{geometry}
\usepackage[table]{xcolor}
\usepackage{fontspec}
\usepackage[spanish]{babel}
\usepackage{tikz}
% Librerías TikZ para los diagramas de las guías (bloque `recursos:`):
%  · babel  → babel en español vuelve activos caracteres como «>», y TikZ los usa
%             en sus opciones (>=stealth, ->). Sin esta librería, cualquier
%             diagrama con flechas revienta con «\language@active@arg> has an
%             extra }». Debe ir ANTES de que se dibuje nada.
%  · positioning        → sintaxis «below left=2pt and 3pt».
%  · decorations.pathreplacing → llaves ({brace}) para señalar rangos.
\usetikzlibrary{babel,positioning,decorations.pathreplacing}
\usepackage{graphicx}
% Circuitos: el trabajo de electrónica se hace hoy con micro:bit y sus sensores
% integrados (MakeCode, sin cableado), así que no se necesitan esquemáticos.
% Si algún día entran montajes con protoboard, basta descomentar esta línea:
% los diagramas .tex/.tikz declarados en `recursos:` ya se incrustan solos.
% \usepackage{circuitikz}
\usepackage[most]{tcolorbox}
\usepackage{tabularx}
\usepackage{array}
\usepackage{fancyhdr}
\usepackage{titlesec}
\usepackage{ragged2e}
\usepackage{enumitem}
\usepackage{microtype}

% Helvetica Neue no trae la flecha U+2192, asi que cada «→» escrito literalmente
% en el YAML se compilaba a NADA, en silencio (462 flechas en 61 guias: rutas del
% tipo «paleta → tipografia → lockup» salian rotas). Mapeamos el caracter a su
% version matematica, que si tiene glifo. Asi el autor puede seguir escribiendo
% «→» comodamente en el YAML.
\usepackage{newunicodechar}
\newunicodechar{→}{\ensuremath{\rightarrow}}

\setmainfont{Helvetica Neue}
\setsansfont{Helvetica Neue}
\setlength{\parindent}{0pt}
\setlength{\parskip}{6pt}
\hyphenpenalty=200
\exhyphenpenalty=200
\pretolerance=400
\tolerance=2000
\setlength{\emergencystretch}{1em}
\renewcommand{\arraystretch}{1.25}
\setlist{leftmargin=5mm,itemsep=1mm,topsep=1mm}
\newcolumntype{Y}{>{\RaggedRight\arraybackslash}X}

\definecolor{milcMagenta}{HTML}{E6007E}
\definecolor{milcVino}{HTML}{5A0038}
\definecolor{milcTurquesa}{HTML}{0093A5}
\definecolor{milcVerde}{HTML}{58A923}
\definecolor{milcMostaza}{HTML}{E5B400}
\definecolor{milcOcre}{HTML}{B86600}
\definecolor{milcNegro}{HTML}{241816}
\definecolor{milcGris}{HTML}{F7F7F4}
\definecolor{milcLinea}{HTML}{E8E2DD}
\definecolor{milcGradePrimary}{HTML}{FF2D87}
\definecolor{milcGradeDark}{HTML}{9F1B56}
\definecolor{milcGradeSoft}{HTML}{FFE0EE}
\definecolor{milcGradeAccent}{HTML}{CFE7FF}
\definecolor{milcGradeText}{HTML}{FFFFFF}
\color{milcNegro}

\pagestyle{fancy}
\fancyhf{}
\lhead{\small Tecnología e Informática · Grado 10}
\rhead{\small Guía 14 · MILC v3}
\cfoot{\small\thepage}
\renewcommand{\headrulewidth}{0pt}

\newtcolorbox{titlebox}[1]{
  enhanced,colback=#1,colframe=#1,arc=7mm,boxrule=0pt,
  left=7mm,right=7mm,top=3mm,bottom=3mm,
  before skip=8pt,after skip=8pt,
  fontupper=\sffamily\bfseries\Large\color{white}
}

\newtcolorbox{softbox}[3]{
  enhanced,colback=#2,colframe=#1,borderline west={4pt}{0pt}{#1},
  arc=2mm,boxrule=.4pt,left=4mm,right=4mm,top=3mm,bottom=3mm,
  before skip=6pt,after skip=8pt,title={#3},coltitle=white,
  colbacktitle=#1,fonttitle=\sffamily\bfseries,fontupper=\RaggedRight,
  attach boxed title to top left={xshift=3mm,yshift=-2mm},
  boxed title style={arc=2mm, boxrule=0pt}
}

\newtcolorbox{drawbox}[2]{
  enhanced,colback=white,colframe=#1,arc=4mm,boxrule=1pt,
  left=5mm,right=5mm,top=4mm,bottom=4mm,before skip=8pt,after skip=8pt,
  title={#2},coltitle=white,colbacktitle=#1,fonttitle=\sffamily\bfseries,
  fontupper=\RaggedRight,
  attach boxed title to top left={xshift=4mm,yshift=-2mm},
  boxed title style={arc=3mm, boxrule=0pt}
}

\newtcolorbox{infoband}[1]{
  enhanced,colback=#1!8,colframe=#1!50,arc=2mm,boxrule=.5pt,
  left=4mm,right=4mm,top=2mm,bottom=2mm,before skip=4pt,after skip=4pt,
  fontupper=\small\RaggedRight
}

% Puente narrativo entre secciones (contrato editorial Conéctate).
% Da continuidad: "Acabas de X. Ahora vas a Y porque Z."
% Visualmente: banda izquierda en color del grado, texto en itálica pequeña.
\newtcolorbox{puentebox}{
  enhanced,colback=milcGradeSoft,colframe=milcGradeDark,
  borderline west={3pt}{0pt}{milcGradeDark},
  arc=0mm,boxrule=0pt,left=5mm,right=4mm,top=2.5mm,bottom=2.5mm,
  before skip=4pt,after skip=8pt,
  fontupper=\itshape\small\color{milcNegro}\RaggedRight
}
% La flecha va en modo matematico a proposito: Helvetica Neue NO tiene el glifo
% U+2192, asi que \textrightarrow se compilaba a NADA («Missing character: There
% is no → in font Helvetica Neue») y el puente salia sin flecha. En modo
% matematico la flecha viene de la fuente matematica y si se dibuja.
\newcommand{\puente}[1]{%
  \begin{puentebox}%
    \textcolor{milcGradeDark}{$\rightarrow$}\hspace{2mm}#1%
  \end{puentebox}%
}

\newcommand{\linead}[1]{\noindent\color{milcLinea}\rule{#1}{0.5pt}\color{milcNegro}}
\newcommand{\respuesta}[1]{\par\vspace{2mm}\foreach \n in {1,...,#1}{\noindent\color{milcLinea}\rule{\linewidth}{0.5pt}\par\vspace{5mm}}\color{milcNegro}}
\newcommand{\checkbox}{\textcolor{milcGradeDark}{\fbox{\phantom{x}}}\hspace{5pt}}

\newcommand{\phasecard}[4]{
\begin{tcolorbox}[enhanced,colback=#3,colframe=#2,arc=3mm,boxrule=.5pt,height=3.05cm,valign=center,halign=center,left=1.5mm,right=1.5mm,top=2mm,bottom=2mm]
{\sffamily\bfseries\fontsize{22}{24}\selectfont\textcolor{#2}{#1}}\par
{\sffamily\bfseries\small #4}
\end{tcolorbox}
}

% ── Recursos visuales de la guía ──────────────────────────────────────
% Declarados en el bloque `recursos:` del YAML y emitidos por el builder.
% \guiaFigura   → imagen rasterizada o PDF (foto, carta celeste, montaje).
% \guiaDiagrama → fuente TikZ/circuitikz (.tex) incrustada como vector:
%                 es la vía para circuitos y esquemáticos editables.
% #1 = ruta relativa al .tex · #2 = pie (puede ir vacío)
\newcommand{\guiaFigura}[2]{%
\begin{center}
\includegraphics[width=0.88\linewidth,height=0.38\textheight,keepaspectratio]{#1}\\[1.5mm]
{\footnotesize\itshape\textcolor{milcNegro!75}{#2}}
\end{center}
}

\newcommand{\guiaDiagrama}[2]{%
\begin{center}
\input{#1}\\[1.5mm]
{\footnotesize\itshape\textcolor{milcNegro!75}{#2}}
\end{center}
}

\begin{document}
\thispagestyle{empty}
\begin{tikzpicture}[remember picture,overlay]
  \fill[white] (current page.south west) rectangle (current page.north east);
  \path[fill=milcGradePrimary,rounded corners=16mm]
    ([xshift=.55cm,yshift=-.55cm]current page.north west) rectangle
    ([xshift=-.55cm,yshift=3.65cm]current page.south east);

  \node[anchor=north west,text=milcGradeText] at ([xshift=2.0cm,yshift=-1.85cm]current page.north west)
    {\sffamily\bfseries\fontsize{14}{16}\selectfont GRADO};
  \node[anchor=north west,text=milcGradeText,opacity=.82] at ([xshift=2.0cm,yshift=-2.75cm]current page.north west)
    {\sffamily\bfseries\fontsize{9}{11}\selectfont SERIE GUÍAS · TECNOLOGÍA E INFORMÁTICA};

  \begin{scope}[shift={([xshift=-3.05cm,yshift=-2.55cm]current page.north east)}]
    \fill[white,opacity=.80] (-.62,-.52) rectangle (.62,.52);
    \draw[milcGradeText,line width=1pt] (-.62,-.52) rectangle (.62,.52);
    \fill[milcGradeDark] (-.42,-.38) rectangle (-.22,.06);
    \fill[milcGradeAccent] (-.10,-.38) rectangle (.10,.24);
    \fill[milcGradePrimary!60!black] (.22,-.38) rectangle (.42,.38);
  \end{scope}

  \node[anchor=west,text=milcGradeText] at ([xshift=1.9cm,yshift=-12.85cm]current page.north west)
    {\sffamily\bfseries\fontsize{124}{124}\selectfont 10};
  % El circulito de grado (°) solo aplica a las guias de grado («11°»); en
  % Bebras y Semillero el numero grande es un momento/modulo, no un grado.
  \draw[milcGradeText,line width=1.7pt]
    ([xshift=4.52cm,yshift=-11.68cm]current page.north west) circle (.13cm);

  \node[anchor=west,text=milcGradeText,text width=8.7cm,align=left] at ([xshift=10.35cm,yshift=-11.6cm]current page.north west)
    {
     {\sffamily\bfseries\fontsize{19}{22}\selectfont Markdown ---\\el formato del oficio\\digital portable}\\[4mm]
     {\sffamily\bfseries\fontsize{23}{26}\selectfont Guía 14-10-TIC}\\[2mm]
     {\sffamily\fontsize{13}{16}\selectfont Período 2 · Informes técnicos y comunicación profesional con IA}\\[5mm]
     {\sffamily\bfseries\fontsize{11}{13}\selectfont Institución Educativa Sor María Juliana}\\[1mm]
     {\sffamily\fontsize{10}{12}\selectfont Autor: Dr. Álvaro Cárdenas Orozco}
    };

  \path[fill=milcGradeSoft,rounded corners=7mm]
    ([xshift=1.35cm,yshift=.85cm]current page.south west) rectangle
    ([xshift=-1.35cm,yshift=4.25cm]current page.south east);
  \draw[milcGradeDark,line width=.65pt,rounded corners=7mm]
    ([xshift=1.35cm,yshift=.85cm]current page.south west) rectangle
    ([xshift=-1.35cm,yshift=4.25cm]current page.south east);
  \node[anchor=center,text=milcNegro,text width=17.0cm,align=center] at ([yshift=2.55cm]current page.south)
    {\sffamily\fontsize{10}{12}\selectfont Metodología MILC · Apertura ancestral · Pensamiento computacional · Triángulo Dussel-Estoicismo-Floridi\\
    \textcolor{milcGradeDark}{\bfseries Producto:} Documento profesional en Markdown con al menos 8 elementos distintos (encabezados de 3 niveles, listas ordenadas y no ordenadas, tabla, enlace, imagen, bloque de código, cita en bloque, énfasis con negrita o cursiva) escrito en un editor gratuito (StackEdit, Dillinger, Obsidian, VS Code) + conversión y exportación a al menos 2 formatos distintos (PDF y HTML) + comparación de 5 líneas con Word/Docs identificando 3 ventajas concretas de Markdown};
\end{tikzpicture}
\mbox{}
\newpage

\section*{Apertura: Markdown --- el formato del oficio digital portable}

\begin{softbox}{milcGradeDark}{milcGradeSoft}{Saber ancestral}
En las sastrerías del centro de Cartago, en los talleres de costura del barrio Obrero, hay un secreto del oficio que separa al sastre profesional del aficionado: \textbf{el uso de patrón estandarizado}. Cuando llega un cliente a hacerse un saco, el sastre no improvisa cortes a ojo. Toma del cajón el \textbf{patrón base} de saco (papel manila grueso, con líneas precisas, calculado matemáticamente). Lo coloca sobre la tela del cliente, ajusta para la talla específica, marca con tiza, corta. El mismo patrón base sirve para 100 clientes distintos: cambia la tela, cambia el ajuste, pero la \textbf{estructura permanece}. Si el sastre intentara coser cada saco improvisando, cada pieza saldría desproporcionada, mal calzada, distinta. El patrón estandarizado da \textbf{portabilidad}: el oficio se puede transmitir, repetir, escalar. Esa práctica ancestral del sastre tiene equivalente en el oficio digital: se llama \textbf{Markdown}. Es el patrón estándar de la escritura digital, que permite que el mismo archivo se convierta a PDF, HTML, libro, presentación, sin perder estructura.
\end{softbox}

\begin{softbox}{milcTurquesa}{milcGris}{Saber contemporáneo}
\textbf{Markdown} es un lenguaje de marcado ligero creado por John Gruber en 2004. Su filosofía: \textbf{escribir texto plano que se convierta automáticamente a formato visual rico}. En lugar de hacer clic en botones para poner \emph{negrita} o \emph{título}, escribes símbolos simples: \texttt{**negrita**}, \texttt{\# Título}, \texttt{- elemento de lista}. El editor o conversor traduce esos símbolos a formato. Su valor profesional es la \textbf{portabilidad}: el mismo archivo \texttt{.md} se convierte a PDF, HTML, libro, presentación, blog. No depende de Word ni de Docs. Sus \textbf{8 elementos básicos} que todo profesional debe conocer: (1) \textbf{Encabezados}: \texttt{\#} para H1, \texttt{\#\#} para H2, \texttt{\#\#\#} para H3. (2) \textbf{Listas no ordenadas}: \texttt{-} o \texttt{*} al inicio de línea. (3) \textbf{Listas ordenadas}: \texttt{1.}, \texttt{2.}, \texttt{3.}. (4) \textbf{Énfasis}: \texttt{*cursiva*} o \texttt{**negrita**}. (5) \textbf{Enlaces}: \texttt{\char91texto\char93(url)}. (6) \textbf{Imágenes}: \texttt{!\char91alt\char93(url)}. (7) \textbf{Bloques de código}: backticks. (8) \textbf{Cita en bloque}: \texttt{>} al inicio. (9) \textbf{Tablas}: con \texttt{|} y guiones. Editores gratuitos: StackEdit (web), Dillinger (web), Obsidian (descargable), VS Code (descargable, gratuito).
\end{softbox}

\begin{softbox}{milcMostaza}{milcGris}{Pregunta-puente}
\textit{¿Qué sabía el sastre al usar patrón estandarizado para clientes distintos, que el novato olvida cuando guarda su escritura solo en formato Word? ¿Y por qué Markdown es el patrón estándar del oficio digital aunque no se vea tan elegante como Word a primera vista?}
\end{softbox}

\begin{softbox}{milcVerde}{milcGris}{Saber y saber hacer}
Al terminar podrás: (1) \textbf{identificar} los 8 elementos básicos de Markdown y cuándo usar cada uno; (2) \textbf{explicar} con tus palabras qué diferencia hay entre escribir en Markdown y en Word/Docs en términos de portabilidad; (3) \textbf{aplicar} la sintaxis Markdown para producir un documento profesional con los 8 elementos; (4) \textbf{evaluar} el documento exportándolo a 2 formatos (PDF y HTML) y comparando con Word/Docs.
\end{softbox}



\small
\begin{tabularx}{\textwidth}{>{\bfseries}p{3.0cm}Y}
\rowcolor{milcGradePrimary!12}Autor & Dr. Álvaro Cárdenas Orozco\\
DBA & Comunicación profesional con asistencia de IA y herramientas ofimáticas gratuitas (MEN, T\&I)\\
Referentes & Markdown como estándar abierto · Google Workspace · Estoicismo (claridad)\\
Producto final & Documento profesional en Markdown con al menos 8 elementos distintos (encabezados de 3 niveles, listas ordenadas y no ordenadas, tabla, enlace, imagen, bloque de código, cita en bloque, énfasis con negrita o cursiva) escrito en un editor gratuito (StackEdit, Dillinger, Obsidian, VS Code) + conversión y exportación a al menos 2 formatos distintos (PDF y HTML) + comparación de 5 líneas con Word/Docs identificando 3 ventajas concretas de Markdown\\
\end{tabularx}
\normalsize

\puente{Antes de empezar, mira el viaje completo. En la sesión 3 trabajaste Google Docs profesional. Hoy aprendes Markdown como alternativa portable. No reemplaza a Docs: lo complementa. Las próximas sesiones (encuestas, análisis con IA, informe comercial, correo profesional) usarán según el caso una u otra.}

\begin{titlebox}{milcGradeDark}
Ruta MILC: así vamos a aprender
\end{titlebox}
\begin{tabularx}{\textwidth}{>{\centering\arraybackslash}X>{\centering\arraybackslash}X>{\centering\arraybackslash}X>{\centering\arraybackslash}X}
\phasecard{1}{milcVerde}{milcGris}{Escucha\\[-1mm]\small Observo, escucho y pregunto.} &
\phasecard{2}{milcTurquesa}{milcGris}{Sistematización\\[-1mm]\small Pensamiento computacional.} &
\phasecard{3}{milcMagenta}{milcGris}{Praxis\\[-1mm]\small Construyo y aplico.} &
\phasecard{4}{milcVino}{milcGris}{Evaluación\\[-1mm]\small Reflexión liberadora.}
\end{tabularx}


\puente{Antes de teorizar, vas a hacer un \emph{experimento simple}: abre StackEdit (\texttt{stackedit.io}, gratis sin cuenta) y en el panel izquierdo escribe: \texttt{\# Mi primer Markdown} en la primera línea, \texttt{\#\# Subtítulo} en la segunda, \texttt{- elemento 1} y \texttt{- elemento 2} debajo. Mira el panel derecho: se renderiza automáticamente. Acabas de hacer Markdown.}

\begin{titlebox}{milcVerde}
Fase 1 · Escucha
\end{titlebox}

\begin{softbox}{milcVerde}{milcGris}{Escena para escuchar}
\textbf{Actividad 1 · IDENTIFICA --- Tu primer Markdown} (15 min · individual con dispositivo). Abre StackEdit en \texttt{stackedit.io} (gratis, sin cuenta). En el panel izquierdo escribe esto exactamente: línea 1: \texttt{\# Mi primer documento}, línea 2: \texttt{Soy estudiante de grado 10}, línea 3: \texttt{\#\# Mis materias favoritas}, línea 4: \texttt{- Tecnología}, línea 5: \texttt{- Literatura}, línea 6: \texttt{**Nota:** este texto está en formato Markdown}. Mira el panel derecho: el texto se renderiza automáticamente. Anota qué símbolo produjo cada efecto visual.
\end{softbox}

\checkbox Abrí StackEdit y escribí 6 líneas en Markdown.\\
\checkbox Vi el renderizado en el panel derecho.\\
\checkbox Identifiqué qué símbolo produjo qué efecto.

\begin{infoband}{milcVerde}


\textbf{Lo que va al cuaderno (Actividad 1 · IDENTIFICA).} Título: «Actividad 1 --- Mi primer Markdown». \textbf{Formato:} captura de los 2 paneles (texto y renderizado) + tabla con símbolos y sus efectos. \textbf{Sabes que terminaste cuando} ves que escribir Markdown es como escribir texto con marcas simples.
\end{infoband}

\puente{Ya hiciste el experimento. Ahora vas a entender la \textbf{anatomía de los 8 elementos básicos} y por qué Markdown es portable.}

\begin{titlebox}{milcTurquesa}
Fase 2 · Sistematización con pensamiento computacional
\end{titlebox}

\begin{softbox}{milcTurquesa}{milcGris}{Cuatro pilares aplicados al tema}
Tu experimento te muestra cómo funciona Markdown. Ahora necesitas la \textbf{anatomía completa} de los 8 elementos profesionales.
\end{softbox}

\small
\begin{tabularx}{\textwidth}{>{\bfseries\RaggedRight\arraybackslash}p{3.6cm}Y}
\rowcolor{milcTurquesa!90}\color{white}Pilar & \color{white}\bfseries Aplicación al tema\\
1. Descomposición & Los \textbf{8 elementos básicos de Markdown} se descomponen así: (1) \textbf{Encabezados}: \texttt{\#} (H1, título grande), \texttt{\#\#} (H2), \texttt{\#\#\#} (H3). (2) \textbf{Énfasis}: \texttt{*cursiva*} o \texttt{\_cursiva\_}; \texttt{**negrita**} o \texttt{\_\_negrita\_\_}. (3) \textbf{Listas no ordenadas}: \texttt{-} o \texttt{*} al inicio de línea, con sangría para subniveles. (4) \textbf{Listas ordenadas}: \texttt{1.}, \texttt{2.}, \texttt{3.} (los números pueden ser todos 1, el renderizador los corrige). (5) \textbf{Enlaces}: \texttt{\char91texto del enlace\char93(url)}. (6) \textbf{Imágenes}: \texttt{!\char91texto alternativo\char93(url o ruta)}. (7) \textbf{Bloques de código}: 3 backticks al abrir y 3 al cerrar, o backtick simple para código en línea. (8) \textbf{Cita en bloque}: \texttt{>} al inicio de cada línea de la cita. Bonus: \textbf{tablas} con pipes \texttt{|} y guiones para separar encabezado.\\
2. Reconocimiento de patrones & La \textbf{portabilidad} de Markdown es su valor único. El mismo archivo \texttt{.md} puede: (a) abrirse en cualquier editor de texto (Bloc de notas, Vim, Sublime, VS Code). (b) renderizarse en GitHub automáticamente. (c) exportarse a PDF con Pandoc o desde StackEdit/Dillinger. (d) convertirse a HTML para publicar en web. (e) usarse en Obsidian para notas conectadas. (f) ser leído por IA con mejor estructura que Word. Word no tiene esta portabilidad: si pierdes Word, pierdes acceso fácil al archivo.\\
3. Abstracción & Las \textbf{3 ventajas concretas} de Markdown frente a Word/Docs: (a) \textbf{Portabilidad universal}: cualquier sistema operativo, cualquier editor, mismo resultado. (b) \textbf{Texto plano legible}: el archivo es lectura humana incluso sin programa especial. (c) \textbf{Versionable con Git}: el control de versiones funciona perfecto con texto plano, no con archivos binarios Word.\\
4. Algoritmo conceptual & Algoritmo: (1) abrir editor (StackEdit es buen primer paso). (2) escribir documento con los 8 elementos. (3) verificar renderizado en panel derecho. (4) exportar a PDF (StackEdit tiene botón). (5) exportar a HTML. (6) abrir los 3 archivos (.md, .pdf, .html) y comparar visualmente. (7) escribir comparación con Word/Docs.\\
\end{tabularx}
\normalsize

\begin{drawbox}{milcTurquesa}{Anatomía de Markdown}
\textbf{Encabezados:} \#, \#\#, \#\#\#. \\[1mm] \textbf{Énfasis:} *cursiva* y **negrita**. \\[1mm] \textbf{Listas:} - / 1. \\[1mm] \textbf{Enlaces:} [texto](url). \\[1mm] \textbf{Imágenes:} ![alt](url). \\[1mm] \textbf{Código:} backticks. \\[1mm] \textbf{Cita:} >. \\[1mm] \textbf{Tabla:} | con guiones.
\end{drawbox}

\begin{softbox}{milcMostaza}{milcGris}{Errores comunes y su corrección}
(1) \textbf{Mezclar Markdown con HTML innecesariamente}: si \texttt{\#} funciona, no uses \texttt{$<$h1$>$}. (2) \textbf{Encabezados sin línea en blanco}: muchos parsers exigen línea vacía antes y después de encabezado. (3) \textbf{Listas sin sangría correcta}: subniveles necesitan 2 o 4 espacios. (4) \textbf{Imágenes con ruta rota}: la ruta relativa o URL debe existir. (5) \textbf{No probar el render}: escribir Markdown sin verificar cómo se ve genera errores que solo se notan al exportar.
\end{softbox}

\begin{infoband}{milcTurquesa}


\textbf{Lo que va al cuaderno (Actividad 2 · EXPLICA).} Título: «Actividad 2 --- Anatomía Markdown». \textbf{Formato:} (a) los 8 elementos con sintaxis; (b) las 3 ventajas; (c) los 5 errores con marca del que más temes. \textbf{Sabes que terminaste cuando} podrías escribir un documento profesional sin abrir esta guía.
\end{infoband}

\puente{Tienes el método. Toca aplicarlo: escribir un documento profesional en Markdown usando los 8 elementos, exportar a PDF y HTML, comparar con Word/Docs.}

\begin{titlebox}{milcMagenta}
Fase 3 · Praxis con pensamiento computacional
\end{titlebox}

\begin{softbox}{milcMagenta}{milcGris}{Construyendo el producto con los cuatro pilares}
\textbf{Actividad 3 · APLICA + EVALÚA --- Documento Markdown + exportación + comparación} (40 min · individual con dispositivo). Hora de aplicar. Escribes documento con los 8 elementos, exportas a PDF y HTML, comparas con Word.
\end{softbox}

\small
\begin{tabularx}{\textwidth}{>{\bfseries\RaggedRight\arraybackslash}p{3.6cm}Y}
\rowcolor{milcMagenta!90}\color{white}Pilar & \color{white}\bfseries Aplicación al producto\\
Descomposición & Tu documento Markdown se construye en 5 pasos: (1) elegir tema (puede ser el informe técnico de la sesión 2, una mini-biografía de alguien que admires, una guía para algo cotidiano). (2) abrir StackEdit o Dillinger. (3) escribir el documento con al menos 8 elementos distintos. (4) exportar a PDF (botón en StackEdit). (5) exportar a HTML.\\
Patrones & Los 8 elementos obligatorios en tu documento: encabezados H1 H2 H3, lista no ordenada, lista ordenada, énfasis bold/cursiva, enlace, imagen (puede ser placeholder como \texttt{!\char91Mi imagen\char93(imagen.png)}), bloque de código, cita en bloque. Bonus: tabla.\\
Abstracción & La comparación con Word/Docs se hace en 5 líneas escritas como cierre: \emph{``Markdown me parece útil para X, X, X. Word/Docs me sigue siendo útil para Y, Y. La principal ventaja de Markdown es Z''}. Esa reflexión te ayuda a entender cuándo usar uno u otro.\\
Algoritmo de armado & Algoritmo: (1) elegir tema. (2) abrir editor. (3) escribir con 8 elementos. (4) verificar render. (5) exportar PDF. (6) exportar HTML. (7) comparar visualmente los 3 archivos (.md, .pdf, .html). (8) escribir comparación con Word.\\
\end{tabularx}
\normalsize

\begin{drawbox}{milcMagenta}{Checklist antes de entregar}
\checkbox Documento Markdown con 8+ elementos distintos. \\[1mm] \checkbox Encabezados en 3 niveles. \\[1mm] \checkbox Listas ordenada y no ordenada. \\[1mm] \checkbox Enlace y imagen. \\[1mm] \checkbox Bloque de código y cita. \\[1mm] \checkbox Exportado a PDF. \\[1mm] \checkbox Exportado a HTML. \\[1mm] \checkbox Comparación con Word de 5 líneas.
\end{drawbox}

\begin{softbox}{milcMagenta}{milcGris}{Plantilla de trabajo}
\textbf{Editor.} \textit{StackEdit, Dillinger, Obsidian, VS Code.} \\[1mm] \textbf{8 elementos.} \textit{Encabezados, listas, énfasis, enlace, imagen, código, cita, tabla.} \\[1mm] \textbf{Exportación.} \textit{PDF + HTML.} \\[1mm] \textbf{Comparación.} \textit{5 líneas Markdown vs Word.}
\end{softbox}

\begin{infoband}{milcMagenta}


\textbf{Lo que va al cuaderno (Actividad 3 · APLICA + EVALÚA).} Título: «Actividad 3 --- Mi documento Markdown». \textbf{Formato:} (a) referencia al archivo .md; (b) captura de PDF y HTML; (c) comparación con Word. \textbf{Sabes que terminaste cuando} entiendes cuándo Markdown te conviene más que Word.
\end{infoband}

\puente{Lo que sigue resume \emph{qué} entregas hoy: documento .md + PDF + HTML + comparación con Word. El criterio principal: que entiendas Markdown como herramienta útil y no como complicación innecesaria.}

\begin{titlebox}{milcMostaza}
Producto de la sesión
\end{titlebox}
\textbf{Documento Markdown + PDF + HTML + comparación con Word}.
\begin{softbox}{milcMostaza}{milcGris}{Criterios de calidad}
(1) 8 elementos distintos en el documento. (2) Exportado a PDF correctamente. (3) Exportado a HTML correctamente. (4) Comparación con Word de 5 líneas. (5) 3 ventajas concretas identificadas.
\end{softbox}

\puente{Antes de cerrar, mira Markdown desde las cinco dimensiones humanas. Aprender un estándar abierto es libertad profesional; depender solo de Word es lock-in con una empresa.}

\begin{titlebox}{milcVino}
Fase 4 · Evaluación liberadora — cinco dimensiones
\end{titlebox}

\begin{softbox}{milcVino}{milcGris}{Sentido de la evaluación}
Evaluar no es solo revisar una nota. Es reconocer qué aprendí, qué cambió en mí, qué emociones manejé, cómo me relaciono con la ciudad y la comunidad, y qué le dejaré a quien viene detrás.
\end{softbox}

\textbf{Mi reflexión liberadora en cinco dimensiones.} Completa cada frase iniciada:

\small
\begin{tabularx}{\textwidth}{>{\bfseries\RaggedRight\arraybackslash}p{3.4cm}Y>{\RaggedRight\arraybackslash}p{4.6cm}}
\rowcolor{milcVino}\color{white}Dimensión & \color{white}\bfseries Mi respuesta (frame starter) & \color{white}\bfseries Algo que descubrí\\
1. Personal & Lo que más cambió en mí fue \linead{6.5cm} & \linead{4.2cm}\\
2. Emocional & La emoción más fuerte fue \linead{2.5cm} y la regulé \linead{3cm} & \linead{4.2cm}\\
3. Ciudadana & En democracia esto importa porque \linead{6cm} & \linead{4.2cm}\\
4. Local (Cartago) & En mi barrio sería útil para \linead{6.5cm} & \linead{4.2cm}\\
5. Intergeneracional & A mi abuela/abuelo le diría \linead{6.5cm} & \linead{4.2cm}\\
\end{tabularx}
\normalsize

\begin{infoband}{milcVino}
\textbf{Para la próxima sustentación.} Vuelve a esta tabla en seis meses. Verás que las respuestas cambian — no porque lo aprendido cambie, sino porque tú cambias. La evaluación liberadora es una conversación contigo a través del tiempo.
\end{infoband}


\puente{Y para terminar, tres voces sobre los estándares abiertos. Dussel sobre las herramientas que liberan al usuario, Séneca sobre la simplicidad como virtud, Floridi sobre los formatos abiertos como ética informacional.}

\begin{titlebox}{milcGradeDark}
Triángulo de pensamiento · cierre filosófico
\end{titlebox}

\begin{softbox}{milcVino}{milcGris}{Dussel · lente del nosotros}
\textit{``Los estándares abiertos liberan al usuario; los formatos cerrados dependen de la empresa que los controla.''}

Word es propiedad de Microsoft; si Microsoft cambia el formato, tus documentos viejos pueden quedar inservibles. Markdown es estándar abierto, no propiedad de nadie: tus archivos seguirán siendo legibles dentro de 50 años con cualquier editor. La filosofía de la liberación pide preferir estándares abiertos cuando es posible. La phronesis del oficio digital consiste en \textbf{no encerrarse en formatos propietarios cuando hay alternativas abiertas}.

\textbf{¿Estoy aprendiendo Markdown como acto de libertad profesional, o lo veo como complicación innecesaria?}
\end{softbox}
\linead{\linewidth}\\[1mm]
\linead{\linewidth}

\begin{softbox}{milcMostaza}{milcGris}{Estoicismo · Séneca · cuidado interior}
\textit{``Lo simple es virtud; la complicación innecesaria es vicio del que prefiere lo difícil por orgullo.''}

Markdown puede parecer demasiado simple comparado con Word. La sabiduría estoica recomienda lo opuesto: \emph{aprecia la simplicidad cuando hace el trabajo}. El sastre con patrón estándar produce trabajo igual de bueno que el que improvisa, pero más rápido y portable. Markdown es ese patrón estándar.

\textbf{¿Estoy apreciando la simplicidad de Markdown, o lo rechazo porque parece demasiado básico?}
\end{softbox}
\linead{\linewidth}\\[1mm]
\linead{\linewidth}

\begin{softbox}{milcTurquesa}{milcGris}{Floridi · lente de la infoesfera}
\textit{``Los formatos abiertos son la nueva ética informacional en la era del control corporativo de datos.''}

En la era de la información, las grandes corporaciones controlan formatos (Word, Photoshop, PowerPoint). La ética digital pide alternativas abiertas que no dependan de una empresa específica. Markdown es ejemplo: gratuito, abierto, portable. Tu aprendizaje de hoy te entrena en una práctica que rige toda profesión técnica seria: \textbf{preferir formatos abiertos para tu información}.

\textbf{¿Estoy contribuyendo a la cultura de formatos abiertos, o sostengo el control corporativo de mis archivos?}
\end{softbox}
\linead{\linewidth}\\[1mm]
\linead{\linewidth}

\begin{titlebox}{milcVino}
Compromiso final
\end{titlebox}

\small
\begin{tabularx}{\textwidth}{>{\RaggedRight\arraybackslash}p{3.0cm}YYY}
\rowcolor{milcVino}\color{white}\bfseries Criterio & \color{white}\bfseries Lo logré & \color{white}\bfseries En proceso & \color{white}\bfseries Necesito apoyo\\
Comprensión del tema & Explico ideas centrales con mis palabras. & Reconozco ideas pero falta precisión. & Necesito repasar conceptos base.\\
Pensamiento computacional & Apliqué los 4 pilares al diseñar mi producto. & Apliqué algunos pilares. & Necesito releer la fase 2.\\
Aplicación y producto & Mi producto cumple los criterios y se puede revisar. & Producto funcional con ajustes pendientes. & Aún en construcción.\\
Reflexión 5 dimensiones & Respondí cada dimensión con honestidad. & Respondí algunas con detalle. & Mis respuestas son muy generales.\\
Triángulo de pensamiento & Conecté las 3 voces con mi trabajo real. & Conecté al menos una. & No relaciono las voces con mi proceso.\\
\end{tabularx}
\normalsize

\textbf{Hoy entendí que:}\\[1mm]
\linead{\linewidth}\\[1mm]
\linead{\linewidth}

\textbf{Mi compromiso para la próxima sesión es:}\\[1mm]
\linead{\linewidth}\\[1mm]
\linead{\linewidth}

\begin{softbox}{milcMostaza}{milcGris}{Cita de despedida · Marco Aurelio}
\textit{``El obstáculo en el camino se vuelve el camino. Lo que estorba al camino se vuelve camino.''}\\[1mm]
Cada error de hoy, bien observado, se convierte en pieza del camino que pisarás mañana.
\end{softbox}

\small
\begin{tabularx}{\textwidth}{>{\bfseries}p{3.6cm}Y}
\rowcolor{milcGradeDark!12}Nombre completo: & \linead{10cm}\\
Fecha: & \linead{4cm} \quad Curso/grupo: \linead{4cm}\\
Mi firma: & \linead{9cm}\\
\end{tabularx}
\normalsize

\begin{infoband}{milcGradeDark}
\textbf{Para la próxima sesión.} Llega con tu documento Markdown, PDF, HTML y comparación. En la sesión 5 vas a aprender a hacer \textbf{encuestas con phronesis} usando Google Forms más IA para detectar preguntas sesgadas. Markdown de hoy es para escribir; Forms de mañana es para preguntar bien.
\end{infoband}

\end{document}
